\documentclass[11pt]{article}
\usepackage[a4paper,margin=23mm]{geometry}
\usepackage{amsmath,amssymb,amsthm,mathtools,bm}
\usepackage{booktabs,array,microtype,xurl,hyperref}
\hypersetup{colorlinks=true,linkcolor=blue,citecolor=blue,urlcolor=blue}

\newtheorem{theorem}{Theorem}[section]
\newtheorem{proposition}[theorem]{Proposition}
\newtheorem{lemma}[theorem]{Lemma}
\newtheorem{corollary}[theorem]{Corollary}
\newtheorem{definition}[theorem]{Definition}
\newtheorem{remark}[theorem]{Remark}
\newcommand{\dd}{\,\mathrm d}
\newcommand{\RR}{\mathbb R}
\newcommand{\SSS}{\mathbb S}
\newcommand{\dist}{\operatorname{dist}}
\newcommand{\osc}{\operatorname{osc}}
\newcommand{\vol}{\operatorname{vol}}

\title{AP-E14 Wess--Zumino Flux Decay and Morrey Regularity:\\
The Exact Hopf Base/Fibre Split and a Div--Curl No-Go}
\author{Route-F derivation ledger}
\date{20 July 2026}

\begin{document}
\maketitle

\begin{abstract}
This note executes the two continuations left by AP-E13.  First, for every
\emph{fixed} complete-minor graph map $u:\Omega\to\SSS^3$, including the
selected relaxed $B=1$ representative, the local Wess--Zumino flux obeys
\[
 \sup_x\frac{\dist(\int_{B_r(x)}u^*\omega_3,
 2\pi^2\mathbb Z)^2}{r^3}
 \leq \frac{4\pi}{3}\sup_x\int_{B_r(x)}|M_3(Du)|^2\dd x=o(1).
\]
This is an unconditional consequence of $L^2$ absolute continuity and does
not use minimality.  It does not imply the linear target oscillation needed
by the AP-E13 replacement: a localized circle-valued rank-one defect has
$M_2=M_3=0$, zero Wess--Zumino flux, finite Dirichlet energy, and oscillation
diameter two at every scale.  Hence the Morrey statement for the actual
relaxed minimizer remains a genuinely variational regularity problem.

Second, writing $n=h(u)$ for the Hopf projection and $a=2u^*\theta$ for the
fibre one-form gives the exact identities
\[
 |Du|^2=\tfrac14(|a|^2+|Dn|^2),\quad
 |M_2(Du)|^2=\tfrac1{16}(|a\wedge Dn|^2+|da|^2),\quad
 |M_3(Du)|^2=\tfrac1{64}|a\wedge da|^2.
\]
Only the signed Chern--Simons integral has exact differential cancellation.
A periodic lift with $da=a\wedge da=0$ but
$\int|M_2(Du)|^2>0$ isolates the surviving vertical--horizontal shear and
disproves a universal reduction of the full cutoff remainder to an exact
div--curl pairing.  A $12/12$ verification card checks the algebra,
counterexamples, exact local Wess--Zumino quadrature, and three independently
relaxed compatible-action $B=1$ grids.  The finite-grid Morrey proxy improves,
but it is not a continuum theorem.  Hessian, determinant, and portal gates
remain closed.
\end{abstract}

\tableofcontents

\section{Scope and conventions}

Let $\Omega\subset\RR^3$ and let $u:\Omega\to\SSS^3\subset\RR^4$ belong to
the complete-minor graph class
\begin{equation}
 u\in W^{1,2},\qquad M_2(Du),M_3(Du)\in L^2.
 \label{eq:graph-class}
\end{equation}
The derivative part of the Route-E continuum action is
\begin{equation}
 W(Du)=\frac12\bigl(|Du|^2+R|M_2(Du)|^2+K|M_3(Du)|^2\bigr),
 \qquad R,K>0.
 \label{eq:W}
\end{equation}
We orient $\SSS^3$ as the boundary of the standard unit ball in $\RR^4$ and
write $\omega_3$ for its volume form, so
$\int_{\SSS^3}\omega_3=2\pi^2$.  For an $\SSS^3$-valued graph map,
\begin{equation}
 u^*\omega_3=J_u\dd x,\qquad
 J_u=\det[u,\partial_1u,\partial_2u,\partial_3u],\qquad
 |J_u|\leq |M_3(Du)|.
 \label{eq:Jbound}
\end{equation}
All statements below distinguish a fixed graph representative from an
arbitrary scale-dependent family.  This distinction is precisely why the
AP-E13 cap sequence is not a counterexample to the next theorem.

\section{Unconditional Wess--Zumino flux decay for a fixed map}

For $\Omega'\Subset\Omega$ and $0<r<\dist(\Omega',\partial\Omega)$ define the
local relative class and its least representative by
\begin{equation}
 \mathcal V_u(x,r)=\left[\int_{B_r(x)}J_u\dd y\right]
 \in\RR/(2\pi^2\mathbb Z),\qquad
 \delta_u(x,r)=\dist\left(\int_{B_r(x)}J_u\dd y,
                         2\pi^2\mathbb Z\right).
 \label{eq:local-flux}
\end{equation}

\begin{theorem}[Uniform fixed-map WZ-flux decay]
Every $u$ satisfying \eqref{eq:graph-class} obeys
\begin{equation}
 \boxed{\quad
 \sup_{x\in\Omega'}\frac{\delta_u(x,r)^2}{r^3}
 \leq \frac{4\pi}{3}
       \sup_{x\in\Omega'}\int_{B_r(x)}|M_3(Du)|^2\dd y
 \xrightarrow[r\downarrow0]{}0.\quad}
 \label{eq:wz-decay}
\end{equation}
In particular this applies to any fixed relaxed $B=1$ minimizer represented
in the complete-minor graph class; neither minimality nor an Euler--Lagrange
equation is used.
\end{theorem}

\begin{proof}
The zero lattice element is an admissible comparison in \eqref{eq:local-flux}.
Using \eqref{eq:Jbound} and Cauchy--Schwarz,
\begin{align}
 \delta_u(x,r)^2
 &\leq\left|\int_{B_r(x)}J_u\dd y\right|^2
 \leq |B_r|\int_{B_r(x)}|M_3(Du)|^2\dd y.       \label{eq:wz-cs}
\end{align}
Since $|B_r|/r^3=4\pi/3$, this gives the inequality.  Put
$f=|M_3(Du)|^2\in L^1(\Omega)$ and
\[
 \alpha_f(t)=\sup\left\{\int_E f\dd y:E\subset\Omega\text{ measurable},
 |E|\leq t\right\}.
\]
Absolute continuity of the Lebesgue integral says $\alpha_f(t)\to0$ as
$t\downarrow0$.  Every $B_r(x)$ has measure $4\pi r^3/3$, uniformly in $x$,
so the right-hand side of \eqref{eq:wz-decay} is at most
$(4\pi/3)\alpha_f(4\pi r^3/3)=o(1)$.
\end{proof}

\begin{remark}[Sharp logical scope]
The conclusion is $\delta_u(x,r)=o(r^{3/2})$, uniformly on compact subsets.
AP-E13 instead chose a different filling candidate $u_r$ at each shrinking
scale and kept a fixed amount of $|M_3|^2$ inside the associated ball.
Absolute continuity for one fixed $L^1$ density gives no uniform estimate
over such a non-equiintegrable family.
\end{remark}

The decay can be arbitrarily slow.  For example the radial majorant
\begin{equation}
 f(y)=\frac{\mathbf1_{|y|<1}}{|y|^3\log^2(e/|y|)}
 \label{eq:slow-density}
\end{equation}
is locally integrable and
$\int_{B_r}f\dd y=4\pi/\log(e/r)$.  Thus no power improvement follows from
$L^2$ control alone.

\section{Why WZ decay does not give linear Morrey oscillation}

The AP-E13 recovery lemma needs a normal-ball trace estimate of the form
\begin{equation}
 \osc_{B_r(x)}u\leq \Lambda r
 \label{eq:linear-morrey}
\end{equation}
or a comparable $L^\infty$ bound.  Graph energy, fixed degree, and
\eqref{eq:wz-decay} cannot imply it.

\begin{proposition}[Rank-one zero-flux obstruction]
There exists an $\SSS^3$-valued finite graph-energy map, localizable in a
vacuum cylinder of a smooth compactly supported $B=1$ map, for which
$M_2=M_3=J_u=0$ almost everywhere but the essential target oscillation has
diameter two in every neighbourhood of the cylinder axis.
\end{proposition}

\begin{proof}
In cylindrical coordinates $(r,\varphi,z)$ take
\begin{equation}
 u(r,z)=\bigl(\cos\ell(r),\sin\ell(r),0,0\bigr),\qquad
 \ell(r)=\log\log(e/r),\qquad 0<r<r_0<1.        \label{eq:rank-one}
\end{equation}
Because $u=\gamma\circ\ell$ for a circle-valued curve $\gamma$, $Du$ has
rank at most one; hence every second and third minor and the WZ density
vanish.  Per unit length of the axis,
\begin{equation}
 \int_{0<r<R}|Du|^2\dd x
 =2\pi\int_0^R\frac{\dd r}{r\log^2(e/r)}
 =\frac{2\pi}{\log(e/R)}\xrightarrow[R\downarrow0]{}0.   \label{eq:tail}
\end{equation}
Nevertheless $\ell(r)\to\infty$ and traverses every phase infinitely often,
so every axis neighbourhood has essential image diameter two.

To localize, replace $\ell$ by a scalar phase $\psi(r,z)$ which equals it in
a smaller cylinder and vanishes near the cylinder boundary.  The composite
$\gamma\circ\psi$ remains rank one.  Insert it where a smooth compactly
supported degree-one map is constant.  Smooth radial truncations of $\psi$
have zero local degree and preserve the exterior degree, so the resulting
graph limit has $B=1$.
\end{proof}

This proposition does \emph{not} prove that the selected relaxed minimizer is
irregular.  It proves the sharper methodological statement that its
minimality must be used.  At present the relaxed functional is not known to
admit the fixed-trace local recovery required to identify it with the naive
integral \eqref{eq:W}; consequently inserting the formal classical
Euler--Lagrange equation here would be circular.

\section{Exact Hopf base/fibre decomposition}

Write $u=(z,w)=(x_1+ix_2,x_3+ix_4)\in\SSS^3\subset\mathbb C^2$ and define
\begin{align}
 h(u)=n&=\bigl(2\Re(z\bar w),\,2\Im(z\bar w),\,|z|^2-|w|^2\bigr),
 \label{eq:hopf}\\
 \theta&=-x_2\dd x_1+x_1\dd x_2-x_4\dd x_3+x_3\dd x_4,
 \qquad a=2u^*\theta.                                  \label{eq:connection}
\end{align}
For the standard orientation
$\omega_2=\frac12\epsilon_{abc}n_a\dd n_b\wedge\dd n_c$ on $\SSS^2$,
direct differentiation gives
\begin{equation}
 da=-n^*\omega_2.                                      \label{eq:curvature}
\end{equation}
Under the fibre-coordinate change $u\mapsto e^{i\chi}u$,
\begin{equation}
 n\mapsto n,\qquad a\mapsto a+2\dd\chi.               \label{eq:gauge}
\end{equation}
This is a useful coordinate split, but not a gauge redundancy of the full
Route-E model: the lift $u$ is physical, and the mass term as well as the
$|a|^2$ energy can depend on its fibre phase.

\begin{theorem}[Complete-minor Hopf identities]
For every tangent derivative $Du\in\RR^{4\times3}$,
\begin{align}
 |Du|^2
 &=\frac14\bigl(|a|^2+|Dn|^2\bigr),                    \label{eq:hopf1}\\
 |M_2(Du)|^2
 &=\frac1{16}\bigl(|a\wedge Dn|^2+|da|^2\bigr),       \label{eq:hopf2}\\
 |M_3(Du)|^2
 &=\frac1{64}|a\wedge da|^2,                          \label{eq:hopf3}\\
 \det[u,Du]
 &=\frac18*(a\wedge da),\qquad
 B=\frac1{16\pi^2}\int_\Omega a\wedge da.            \label{eq:hopfdegree}
\end{align}
The sign in \eqref{eq:hopfdegree} follows from the stated $\theta$ and the
boundary orientation of $\SSS^3$; reversing either orientation reverses both
displayed signs and changes no energy identity.
\end{theorem}

\begin{proof}
The identities are invariant under the transitive $U(2)$ action, so evaluate
at $u=(1,0,0,0)$.  A tangent column has coordinates
$p=(0,v,p_3,p_4)$.  At this point
\begin{equation}
 a(p)=2v,\qquad Dn(p)=(2p_3,-2p_4,0),\qquad
 da(p,q)=4(p_3q_4-p_4q_3).                              \label{eq:frame}
\end{equation}
The vertical line and horizontal two-plane are orthogonal, proving
\eqref{eq:hopf1}.  For each pair of domain columns, the squared two-area is
the sum of the vertical--horizontal parallelogram and horizontal determinant
squares.  Equation \eqref{eq:frame} multiplies each by sixteen, proving
\eqref{eq:hopf2}.  The three-volume is the determinant of the one vertical
and two horizontal components; $a\wedge da$ multiplies it by eight.  This
proves \eqref{eq:hopf3} and \eqref{eq:hopfdegree}.
\end{proof}

Consequently \eqref{eq:W} becomes
\begin{equation}
 W=\frac18(|a|^2+|Dn|^2)
  +\frac{R}{32}(|a\wedge Dn|^2+|da|^2)
  +\frac{K}{128}|a\wedge da|^2.                        \label{eq:splitW}
\end{equation}
The split is compatible with Sobolev Hopf-lifting theory
\cite{AucklyKapitanski2008,GuerraLamyZemas2026Lifting} and with recent
special Faddeev--Skyrme minimality results \cite{GuerraLamyZemas2026Hopf},
but those theorems do not supply $L^2$ control of every complete minor in the
present $\SSS^3$ problem.

\section{The exact div--curl test}

The signed Chern--Simons density has the expected exact cancellation:
\begin{equation}
 (a+2\dd\chi)\wedge d(a+2\dd\chi)
 =a\wedge da+2\dd(\chi\,da).                           \label{eq:cs-exact}
\end{equation}
Thus $\int a\wedge da$ is invariant on a torus or under fixed boundary data.
The positive terms in \eqref{eq:splitW} do not share this cancellation.

\begin{proposition}[Chern--Simons exactness does not remove the cutoff stress]
The AP-E13 positive cutoff remainder cannot, in general, be rewritten only
as an exact divergence or a curl-controlled pairing in the Hopf variables.
\end{proposition}

\begin{proof}
Two independent periodic tests isolate the failure.

First let $T^3=(\RR/2\pi\mathbb Z)^3$ and
\begin{equation}
 u=\left(e^{i\chi}\cos\frac t2,
          e^{i\chi}\sin\frac t2\right),\qquad
 t=A\sin x_1,\quad \chi=B\sin x_2.                    \label{eq:shear-lift}
\end{equation}
Then $n=(\sin t,0,\cos t)$ has rank one,
$a=2B\cos x_2\dd x_2$, and
\begin{equation}
 da=a\wedge da=0,\qquad
 \int_{T^3}|M_2(Du)|^2\dd x
 =\frac1{16}\int_{T^3}|a\wedge Dn|^2\dd x
 =\frac{A^2B^2\pi^3}{2}>0.                            \label{eq:shear-positive}
\end{equation}
Every expression depending only on $da$ or $a\wedge da$ vanishes, whereas
the vertical--horizontal shear remains strictly positive.

Second take
\begin{equation}
 a=\sin z\dd x+\cos z\dd y+c\dd z,\qquad
 \chi=\varepsilon\sin x.                              \label{eq:cs-test}
\end{equation}
Here $*(a\wedge da)=1$, while
$*((a+2\dd\chi)\wedge da)=1+2\varepsilon\cos x\sin z$.
The signed integrals both equal $(2\pi)^3$, but
\begin{equation}
 \int_{T^3}|(a+2\dd\chi)\wedge da|^2
 =(2\pi)^3(1+\varepsilon^2)>(2\pi)^3.                 \label{eq:cs-square}
\end{equation}
Thus exactness controls the signed integral, not its squared density.
\end{proof}

In particular the problematic AP-E13 term transforms exactly as
\begin{equation}
 \frac{|u-q|^2}{r^2}|M_2(Du)|^2
 =\frac{|u-q|^2}{16r^2}
   \bigl(|a\wedge Dn|^2+|da|^2\bigr).                 \label{eq:remainder}
\end{equation}
The counterexample disproves the \emph{universal full cancellation}; it does
not exclude a model-specific compensation theorem which also uses relaxed
minimality, a Noether identity, or elimination of a defect measure.

\section{Same-action finite-grid audit}

The verifier re-relaxes the AP-E12 compatible-cochain action, rather than
recycling stationary files.  The normalized-affine map on a tetrahedron $T$
is $u=v/|v|$, $v(\lambda)=\sum_{i=0}^3\lambda_i u_i$.  With orientation
$\sigma_T$, determinant $D_T=\det[u_0,u_1,u_2,u_3]$, and physical volume
$V_T$, the exact radial identities used in quadrature are
\begin{align}
 \int_Tu^*\omega_3
 &=\sigma_TD_T\int_\Delta|v(\lambda)|^{-4}\dd\lambda, \label{eq:tetwz}\\
 \int_T|M_3(Du)|^2\dd x
 &=\frac{D_T^2}{6V_T}\int_\Delta|v(\lambda)|^{-8}\dd\lambda. \label{eq:tetm3}
\end{align}
They reproduce the local Cauchy bound
$|\int_Au^*\omega_3|^2\leq |A|\int_A|M_3|^2$.

\begin{table}[ht]
\centering
\small
\begin{tabular}{@{}rrrrrr@{}}
\toprule
$N$ & $a$ & projected stationarity & $\max\angle/a$ & q99 $\angle/a$
& $\max$ WZ ratio\\
\midrule
17 & 0.343750 & $1.539\times10^{-7}$ & 3.7669 & 1.9147 & 0.8512\\
21 & 0.300000 & $5.003\times10^{-8}$ & 3.4636 & 1.8115 & 0.9481\\
25 & 0.270833 & $3.957\times10^{-8}$ & 3.3162 & 1.7043 & 0.9723\\
\bottomrule
\end{tabular}
\caption{All three backgrounds are independently relaxed, admissible, and
have the three normalized-affine degree evaluations $[1,1,1]$.  The last
column is the largest local ratio $\Phi^2/(|A|\int_A|M_3|^2)$ among the
reported $2a$ and $3a$ balls.}
\label{tab:grid}
\end{table}

At centre-node radii $r=a,2a,3a$, the maximum target-angle quotients are
$(2.202,2.054,1.915)$, $(1.964,1.912,1.849)$, and
$(1.878,1.846,1.814)$ for $N=17,21,25$, respectively.  At $r=2a$ the measured
$\Phi^2/r^3$ decreases $260.31\to150.09\to101.47$.  These trends support a
regular continuum representative, but a three-grid bounded edge quotient is
not a uniform continuum Lipschitz estimate and cannot promote
\eqref{eq:linear-morrey}.

The production card evaluates 128 random tangent jets.  Its largest residual
is $8.53\times10^{-14}$ in \eqref{eq:hopf2}; the largest local WZ Cauchy ratio
is $0.972327697$.  All twelve algebraic, analytic, topological, numerical,
and fail-closed gate checks pass.

\section{Gate decision and the next theorem}

The results separate three claims which AP-E13 had left coupled:
\begin{enumerate}
 \item Fixed-representative WZ-flux decay is \textbf{proved} and needs no
       variational input.
 \item Linear-scale Morrey oscillation for the selected relaxed minimizer is
       \textbf{open}; graph bounds, degree, and WZ decay alone are
       insufficient.
 \item A universal Hopf exact-div--curl reduction of the full positive cutoff
       remainder is \textbf{disproved}.  Only the signed Chern--Simons part is
       exact.
\end{enumerate}

The blocker-changing continuation is therefore not another reformulation of
the WZ flux.  It is a theorem about the lower-semicontinuous relaxation
itself:
\begin{enumerate}
 \item derive an inner-variation or monotonicity identity for the relaxed
       energy measure, allowing an explicit stress/defect measure;
 \item prove that its rank-one vertical--horizontal Hopf defect vanishes, or
       construct a degree-preserving local comparison which removes it;
 \item obtain an $\varepsilon$-regularity decay estimate and then upgrade
       mean Morrey decay to the $L^\infty$ oscillation in
       \eqref{eq:linear-morrey}.
\end{enumerate}
Only after local recovery, classicality, and continuum isolation are proved
may the same-action Riemann Hessian and determinant variation be assembled.
The degree-one portal remains later still.  Recent general Sobolev screening
results \cite{BousquetPonceVanSchaftingen2025} do not presently control the
$L^2$ complete-minor endpoint required here.

\section*{Reproducibility}

Run
\begin{verbatim}
python3 route_f/code/verify_ap_e14_wz_morrey_hopf_divcurl.py
\end{verbatim}
from the repository root.  The machine-readable outputs are
\begin{itemize}
 \item \path{route_f/output/ap_e14_wz_morrey_hopf_divcurl.json},
 \item \path{route_f/output/ap_e14_wz_morrey_hopf_divcurl.md}.
\end{itemize}
Production reports
\path{12/12; wz=True; morrey=False; divcurl=False; hessian=False}.

\bibliographystyle{plain}
\bibliography{route_f/tex/ap_e14_wz_morrey_hopf_divcurl}

\end{document}
